\section{Objective under-determination of skeletal position}
\label{sec:skeleton-underdetermined}

The coxa result in \S\ref{sec:articulation} was pursued to a mechanism
rather than left as an observation; doing so honestly first required
checking the ground truth being pursued against. An early set of expert
joint annotations was recorded in Blender's internal coordinate convention,
one axis permutation, $(x,y,z)\mapsto(x,z,-y)$, away from the convention
used by the scan meshes those annotations were meant to describe; a
derived skeleton reference file inherited the same mismatch. Every
conclusion drawn against that reference in the project's first analysis
pass, including an apparent $42.7\%$ production skeleton error, was
invalidated once this was found, and discarded. This is reported as a
technical result in its own right (\F{F7}), not a caveat: catching a
ground-truth defect before it propagates into a project's headline claims
is exactly the discipline a registration pipeline needs.

\begin{investigation}{I5}{Does the objective actually constrain skeleton
position, once the ground truth used to check it is trustworthy?}
\invQuestion{Is there a solution, reachable within the model's own shape
budget, that is substantially more anatomically accurate than production's,
and if so, why does the optimiser not find it?}
\invMethod{The corrected \term{Joint Alignment Benchmark}: eleven
expert-annotated specimens, re-registered against the mesh coordinate
convention to a residual of $7\times10^{-8}$, frozen and hash-verified,
yielding 396 fits under a fixed protocol. Expert joint targets are supplied
directly during fitting and compared against the unsupervised production
fit; a systematic single-factor sweep re-tests every existing objective term
in isolation; a leave-one-region-out variant withholds expert targets for
one anatomical region at a time.}
\invResult{Production skeleton error is $25.1\%$ of Weber's length (95\% CI
$17.5$--$44.8\%$). Supplying expert joint targets during fitting reduces
this to $17.4\%$ ($10$ of $11$ specimens improved). The single-factor sweep
found no individual existing objective term that reliably improved joint
alignment on its own (Friedman test, $p=.011$, with the unmodified
production configuration retaining the best mean rank across variants).
Leave-one-region-out supervision does not propagate: for every held-out
region, the confidence interval for improvement included zero.}
\invVerdict{REJECT}{The production objective contains no explicit
observation term for joint location. Its surface terms constrain the
skeleton only indirectly, through the skinned surface it produces, and do
not distinguish all anatomically different latent configurations that yield
similar surfaces. No single existing term, reweighted or removed, fixes this
reliably, and supervision supplied for one region does not propagate to
another. Recorded as \F{F6}; it also explains why reweighting existing terms
(\S\ref{sec:eliminated}) returned null results.}
\end{investigation}

The error this benchmark measures is itself structured rather than uniform.
Proximal joints are recovered far more reliably than distal ones:
coxae at $11.6\%$ of Weber's length, trochanter/femur at $12.9\%$, while
the body axis, mandibles, antennae, and tibia/tarsus segments range from
$24\%$ to over $56\%$ (Table~\ref{tab:regional-error}). A log-scale variance
decomposition attributes far more of the spread in this benchmark to
\emph{which joint} and \emph{which specimen} than to \emph{which random
seed} the optimiser started from.

\begin{table}[!htb]
\centering
\caption{Regional skeleton error (\% of Weber's length) on the corrected
benchmark, and the log-scale variance decomposition across the same fits.}
\label{tab:regional-error}
\begin{tabular}{@{}lc@{\hspace{2.5em}}lc@{}}
\toprule
\textbf{Anatomical region} & \textbf{Error} & \textbf{Variance source} & \textbf{Share} \\
\midrule
Coxae               & 11.6\% & Which joint          & 0.36 \\
Trochanter / femur   & 12.9\% & Which specimen       & 0.17 \\
Body axis            & 24.1\% & Optimiser random seed & 0.02 \\
Mandibles            & 42.8\% & & \\
Antennae             & 54.9\% & & \\
Tibia / tarsus       & 56.3\% & & \\
\bottomrule
\end{tabular}
\end{table}

Registration failure is therefore anatomically structured, not attributable
to optimiser luck (\F{F8}): a fit that happens to land badly is far more
likely to have landed on a hard joint or a hard specimen than to have simply
been unlucky in its random initialisation.
